% Options for packages loaded elsewhere
\PassOptionsToPackage{unicode}{hyperref}
\PassOptionsToPackage{hyphens}{url}
\documentclass[
  12pt,
]{ctexart}
\usepackage{xcolor}
\usepackage{amsmath,amssymb}
\setcounter{secnumdepth}{-\maxdimen} % remove section numbering
\usepackage{iftex}
\ifPDFTeX
  \usepackage[T1]{fontenc}
  \usepackage[utf8]{inputenc}
  \usepackage{textcomp} % provide euro and other symbols
\else % if luatex or xetex
  \usepackage{unicode-math} % this also loads fontspec
  \defaultfontfeatures{Scale=MatchLowercase}
  \defaultfontfeatures[\rmfamily]{Ligatures=TeX,Scale=1}
\fi
\usepackage{lmodern}
\ifPDFTeX\else
  % xetex/luatex font selection
\fi
% Use upquote if available, for straight quotes in verbatim environments
\IfFileExists{upquote.sty}{\usepackage{upquote}}{}
\IfFileExists{microtype.sty}{% use microtype if available
  \usepackage[]{microtype}
  \UseMicrotypeSet[protrusion]{basicmath} % disable protrusion for tt fonts
}{}
\makeatletter
\@ifundefined{KOMAClassName}{% if non-KOMA class
  \IfFileExists{parskip.sty}{%
    \usepackage{parskip}
  }{% else
    \setlength{\parindent}{0pt}
    \setlength{\parskip}{6pt plus 2pt minus 1pt}}
}{% if KOMA class
  \KOMAoptions{parskip=half}}
\makeatother
\usepackage{longtable,booktabs,array}
\usepackage{calc} % for calculating minipage widths
% Correct order of tables after \paragraph or \subparagraph
\usepackage{etoolbox}
\makeatletter
\patchcmd\longtable{\par}{\if@noskipsec\mbox{}\fi\par}{}{}
\makeatother
% Allow footnotes in longtable head/foot
\IfFileExists{footnotehyper.sty}{\usepackage{footnotehyper}}{\usepackage{footnote}}
\makesavenoteenv{longtable}
\setlength{\emergencystretch}{3em} % prevent overfull lines
\providecommand{\tightlist}{%
  \setlength{\itemsep}{0pt}\setlength{\parskip}{0pt}}
\usepackage{bookmark}
\IfFileExists{xurl.sty}{\usepackage{xurl}}{} % add URL line breaks if available
\urlstyle{same}
\hypersetup{
  hidelinks,
  pdfcreator={LaTeX via pandoc}}

\author{SCX}
\date{2026}

\begin{document}

\section{SCX
项目成果归属说明}\label{scx-ux9879ux76eeux6210ux679cux5f52ux5c5eux8bf4ux660e}

\begin{quote}
本文件记录 SCX (State-Conditioned eXpertise)
项目的知识产权归属与独立性声明。 最后更新：2026-06-26
\end{quote}

\begin{center}\rule{0.5\linewidth}{0.5pt}\end{center}

\subsection{一、SCX
项目独立性声明}\label{ux4e00scx-ux9879ux76eeux72ecux7acbux6027ux58f0ux660e}

\subsubsection{1.1 核心性质：SCX 是纯数学/ML
框架}\label{ux6838ux5fc3ux6027ux8d28scx-ux662fux7eafux6570ux5b66ml-ux6846ux67b6}

SCX 是\textbf{状态条件专家性}（State-Conditioned
eXpertise）的数学框架，其核心贡献是：

\begin{itemize}
\tightlist
\item
  一个定义体系（SCX Reliability, SCX Data Four-Classification,
  State-Wise Active Learning, Expert Routing, Noise vs Hard State 区分）
\item
  一组形式化命题（3 个核心定理 + 2 个扩展定理）
\item
  对应的 Python 实现
\end{itemize}

SCX \textbf{不包含}任何 DFT（密度泛函理论）计算代码、VASP
输入文件、或材料科学特定的计算逻辑。

\subsubsection{1.2
不依赖学校资源}\label{ux4e0dux4f9dux8d56ux5b66ux6821ux8d44ux6e90}

\begin{longtable}[]{@{}
  >{\raggedright\arraybackslash}p{(\linewidth - 4\tabcolsep) * \real{0.2247}}
  >{\raggedright\arraybackslash}p{(\linewidth - 4\tabcolsep) * \real{0.1348}}
  >{\raggedright\arraybackslash}p{(\linewidth - 4\tabcolsep) * \real{0.6404}}@{}}
\toprule\noalign{}
\begin{minipage}[b]{\linewidth}\raggedright
资源
\end{minipage} & \begin{minipage}[b]{\linewidth}\raggedright
是否使用
\end{minipage} & \begin{minipage}[b]{\linewidth}\raggedright
说明
\end{minipage} \\
\midrule\noalign{}
\endhead
\bottomrule\noalign{}
\endlastfoot
孝感超算（学校设备） & \textbf{否} & SCX
所有理论工作、编码、实验均在个人电脑完成 \\
学校网络/存储 & \textbf{否} & SCX 项目存储于个人电脑和用户个人 GitHub
账户 \\
学校软件许可 & \textbf{否} & SCX 仅使用开源 Python 包（numpy, scipy,
scikit-learn 等） \\
学校人员/学生 & \textbf{否} & SCX 为单人（用户）独立完成 \\
\end{longtable}

\subsubsection{1.3 实验独立性}\label{ux5b9eux9a8cux72ecux7acbux6027}

SCX 的验证实验包括三类：

\begin{enumerate}
\def\labelenumi{\arabic{enumi}.}
\tightlist
\item
  \textbf{合成 2D 数据实验}：数据由
  \texttt{experiments/synthetic/data\_generator.py}
  自行生成，不需要任何外部数据源
\item
  \textbf{MedMNIST / CIFAR}：公开学术基准数据集，任何人可下载使用
\item
  \textbf{MLIP 案例}：使用 AlN v3 DFT 数据（EGP
  项目成果）作为可选的、非必须的附加展示
\end{enumerate}

\begin{quote}
合成数据 + 公开基准数据集已足以验证 SCX 的所有核心主张。MLIP
案例纯粹是''锦上添花''的附加展示。
\end{quote}

\subsubsection{1.4 AI
代码生成声明}\label{ai-ux4ee3ux7801ux751fux6210ux58f0ux660e}

SCX Python 包的\textbf{全部代码}由 AI 工具 \textbf{Claude Code
(Anthropic)} 在用户的数学框架和设计方向指导下生成。用户是：

\begin{itemize}
\tightlist
\item
  \textbf{数学框架的原创者}（定义体系、命题、证明、算法设计）
\item
  \textbf{架构设计者}（模块划分、接口设计、整体架构）
\item
  \textbf{代码的委托生成者}（向 AI 提供设计规格，AI 生成实现）
\end{itemize}

\begin{center}\rule{0.5\linewidth}{0.5pt}\end{center}

\subsection{二、与 EGP / AlGaN / DFT
项目的边界}\label{ux4e8cux4e0e-egp-algan-dft-ux9879ux76eeux7684ux8fb9ux754c}

\subsubsection{2.1
两个独立项目}\label{ux4e24ux4e2aux72ecux7acbux9879ux76ee}

\begin{longtable}[]{@{}lll@{}}
\toprule\noalign{}
维度 & EGP 项目（Paper 1-3） & SCX 项目（Paper 4） \\
\midrule\noalign{}
\endhead
\bottomrule\noalign{}
\endlastfoot
本质 & MLIP（机器学习势函数） & 通用的 ML 理论框架 \\
依赖 & 依赖 DFT 数据（AlN v3） & 不依赖任何 DFT 数据 \\
计算资源 & 使用孝感超算 & 个人电脑即可 \\
应用领域 & 材料科学 & 通用机器学习 \\
代码仓库 & \texttt{egp/} & \texttt{SCX/} \\
\end{longtable}

\subsubsection{2.2
唯一交集：概念起源}\label{ux552fux4e00ux4ea4ux96c6ux6982ux5ff5ux8d77ux6e90}

SCX 中\textbf{``状态条件专家性''}这一核心概念的思想萌芽来自 EGP Paper 1
的 gauge fixing 工作------在拼接多个 ACE/PACE
势函数时，用户观察到同一势函数在不同构型区域上的精度不同。

然而，\textbf{SCX
框架本身完全不涉及势函数、DFT、或任何材料科学内容}。它是一个通用 ML
理论框架，将''状态条件专家性''形式化为可用于任何监督学习场景的数学框架（数据价值评估、主动学习、专家路由、数据压缩）。

\begin{quote}
\textbf{类比}：牛顿在观察苹果落地时萌生了万有引力的想法，但万有引力定律是一个普适物理理论，不依赖于苹果的具体种类。
\end{quote}

\subsubsection{2.3 EGP 项目的 DFT
数据权属}\label{egp-ux9879ux76eeux7684-dft-ux6570ux636eux6743ux5c5e}

AlN v3 DFT 数据是在\textbf{孝感超算}（学校设备）上使用 VASP 生成的，属于
EGP 项目的成果。SCX 项目\textbf{不声称对这部分数据拥有独立所有权}。SCX
框架的可验证性不依赖这些数据。

\begin{center}\rule{0.5\linewidth}{0.5pt}\end{center}

\subsection{三、与学校/课题组的关系说明}\label{ux4e09ux4e0eux5b66ux6821ux8bfeux9898ux7ec4ux7684ux5173ux7cfbux8bf4ux660e}

\subsubsection{3.1 切割声明}\label{ux5207ux5272ux58f0ux660e}

\begin{itemize}
\tightlist
\item
  SCX 项目\textbf{不是}用户所在学校或课题组的资助项目
\item
  SCX 项目\textbf{未使用}学校的研究经费
\item
  SCX 项目\textbf{未使用}学校的计算资源、软件许可或数据
\item
  SCX 项目\textbf{不是}用户学位论文或职务研究的一部分
\item
  SCX 项目\textbf{完全在个人时间、使用个人设备、个人计算资源}完成
\end{itemize}

\subsubsection{3.2
潜在关联与排除}\label{ux6f5cux5728ux5173ux8054ux4e0eux6392ux9664}

SCX
项目与用户所在课题组之间\textbf{没有雇佣、资助或指导关系}。即使存在任何课题组的通用知识（如
ML 基础），SCX
的理论创新------状态条件专家性的数学框架------是用户在个人时间独立提出的原创贡献。

\begin{center}\rule{0.5\linewidth}{0.5pt}\end{center}

\subsection{四、权属声明模板}\label{ux56dbux6743ux5c5eux58f0ux660eux6a21ux677f}

\subsubsection{4.1 所有权归属}\label{ux6240ux6709ux6743ux5f52ux5c5e}

SCX 项目的所有权利（包括但不限于版权、专利权、署名权）归属：

\begin{quote}
\textbf{【预留：填写权利人信息】}

姓名：SCX
日期：\_\_\_\_\_\_\_\_\_\_\_\_\_2026.06.26\_\_\_\_\_\_\_\_\_\_\_\_\_
\end{quote}

\subsubsection{4.2 代码许可}\label{ux4ee3ux7801ux8bb8ux53ef}

SCX Python 包的发布许可待定。目前为私密仓库，不对外分发。

\subsubsection{4.3 论文署名}\label{ux8bbaux6587ux7f72ux540d}

如发表学术论文，SCX 的署名权根据独立贡献程度决定：

\begin{longtable}[]{@{}ll@{}}
\toprule\noalign{}
贡献 & 角色 \\
\midrule\noalign{}
\endhead
\bottomrule\noalign{}
\endlastfoot
数学框架（定义、命题、证明） & \textbf{用户}（100\%） \\
架构设计 & \textbf{用户}（100\%） \\
代码实现 & AI (Claude Code) 按用户指导生成 \\
实验设计和执行 & \textbf{用户}（100\%） \\
\end{longtable}

\begin{quote}
\textbf{AI 工具不享有署名权}。根据学术出版规范（COPE, Nature, arXiv
等），AI 工具不得列为作者，但应在致谢或方法部分声明其使用。
\end{quote}

\begin{center}\rule{0.5\linewidth}{0.5pt}\end{center}

\subsection{五、法律与伦理声明}\label{ux4e94ux6cd5ux5f8bux4e0eux4f26ux7406ux58f0ux660e}

本文件所陈述的内容如实反映 SCX
项目的开发过程、资源使用情况和贡献归属。如有争议，保留通过以下方式进一步证明的权利：

\begin{enumerate}
\def\labelenumi{\arabic{enumi}.}
\tightlist
\item
  Git 提交历史（所有提交来自个人电脑）
\item
  Claude Code 对话记录（AI 生成代码的证据）
\item
  个人电脑文件元数据（创建/修改时间戳）
\item
  开发日志（本仓库中的 \texttt{DEVELOPMENT\_LOG.md}）
\end{enumerate}

\begin{center}\rule{0.5\linewidth}{0.5pt}\end{center}

\emph{本文件由用户于 2026-06-26 撰写，用于记录 SCX 项目的知识产权状况。}

\end{document}
